\documentclass[../../../include/open-logic-section]{subfiles}

\begin{document}

\olfileid{sth}{ord-arithmetic}{using-addition}
\olsection{ક્રમવાચક સરવાળાનો ઉપયોગ}

ક્રમવાચક સંખ્યાઓના સરવાળાથી
\olref[spine][rank]{defnsetrank}ના અર્થમાં જુદા જુદા
ગણોના સ્તરાંક સ્પષ્ટ રીતે ગણી શકીએ:

\begin{lem}\ollabel{rankcomputation}
જો $\setrank{A} = \alpha$ અને $\setrank{B} = \beta$ હોય, તો:
\begin{enumerate}
	\item\ollabel{exrankpow} $\setrank{\Pow{A}} = \alpha\ordplus 1$
	\item\ollabel{exrankpair} $\setrank{\{A, B\}} = \max(\alpha,
	\beta) \ordplus 1$
	\item\ollabel{exrankcup} $\setrank{A \cup B} = \max(\alpha,
	\beta)$
	\item\ollabel{exranktuple} $\setrank{\tuple{A,B}} = \max(\alpha,
	\beta) \ordplus  2$
	\item\ollabel{exranktimes} $\setrank{A \times B} \leq \max(\alpha,
	\beta) \ordplus  2$
	\item\ollabel{exrankunion} $\alpha$ ખાલી અથવા સીમા
	હોય ત્યારે $\setrank{\bigcup A} = \alpha$;
	$\alpha = \gamma\ordplus 1$ હોય ત્યારે
	$\setrank{\bigcup A} = \gamma$
\end{enumerate}
\end{lem}

\begin{proof}
આખી સાબિતીમાં \olref[spine][rank]{ranksupstrict}નો
વારંવાર ઉપયોગ કરીએ છીએ.

\emph{\olref{exrankpow}.} જો $x \subseteq A$ હોય, તો
$\setrank{x} \leq
\setrank{A}$. તેથી $\setrank{\Pow{A}} \leq \alpha \ordplus  1$.
ખાસ કરીને $A
\in \Pow{A}$ હોવાથી $\setrank{\Pow{A}} = \alpha \ordplus  1$.

\emph{\olref{exrankpair}.} \olref[spine][rank]{ranksupstrict} મુજબ.

\emph{\olref{exrankcup}.} \olref[spine][rank]{ranksupstrict} મુજબ.

\emph{\olref{exranktuple}.} \olref{exrankpair} બે વાર વાપરો.

\emph{\olref{exranktimes}.} નોંધો કે
$A \times B \subseteq
\Pow{\Pow{A \cup B}}$; પછી \olref{exranktuple} વાપરો.

\emph{\olref{exrankunion}.} જો $\alpha = \gamma\ordplus 1$
હોય, તો કોઈ $c \in A$ માટે $\setrank{c} = \gamma$
અને $A$ના કોઈ !!{element}નો સ્તરાંક તેનાથી મોટો નથી;
તેથી $\setrank{\bigcup A} = \gamma$. જો $\alpha$
સીમા ક્રમવાચક સંખ્યા હોય, તો $A$માં $\alpha$ની
ગમે તેટલી નજીક, પણ તેનાથી કડક રીતે નાના, સ્તરાંક
ધરાવતા !!{element}s છે. તેથી $\bigcup A$માં પણ
$\alpha$ની ગમે તેટલી નજીક, પણ તેનાથી કડક રીતે
નાના, સ્તરાંક ધરાવતા !!{element}s છે; આથી
$\setrank{\bigcup A} = \alpha$.
\end{proof}
\noindent
\olref{exranktimes}માં સમાનતાને બદલે \emph{અસમાનતા}
શા માટે છે તે સ્વાધ્યાય તરીકે છોડીએ છીએ.

\begin{prob}
$\setrank{A \times B}=
\max(\setrank{A}, \setrank{B})$ થાય એવા ગણો
$A$ અને $B$ શોધો. પછી
$\setrank{A \times B}=\max(\setrank{A}, \setrank{B}) \ordplus  2$
થાય એવા ગણો $A$ અને $B$ શોધો. શું બીજા કોઈ
સ્તરાંકો શક્ય છે?
\sourcecorrection{OLMD-131}{મૂળના બીજા પ્રશ્નમાં
$\setrank{A\times B}$ અને $\max(\setrank{A},\setrank{B})\ordplus2$
વચ્ચે સંબંધનું ચિહ્ન છૂટી ગયું છે. પહેલો પ્રશ્ન અને
ઉપરની ઉચ્ચસીમા-અસમાનતા પ્રમાણે અહીં ``$=$''
ઉમેર્યું છે.}
\end{prob}

હવે ક્રમવાચક સંખ્યાને ``સીમિત'' કે ``અનંત'' કહેવાનો
શું અર્થ થાય તેના ઘણા વાજબી ખ્યાલો સમકક્ષ છે તે
બતાવી શકીએ:

\begin{lem}\ollabel{ordinfinitycharacter}
કોઈ પણ ક્રમવાચક સંખ્યા $\alpha$ માટે નીચેની શરતો
સમકક્ષ છે:
\begin{enumerate}
	\item\ollabel{ord:notinomega} $\alpha\notin \omega$;
	એટલે $\alpha$ પ્રાકૃતિક સંખ્યા નથી
	\item\ollabel{ord:omegaplus} $\omega \leq \alpha$
	\item\ollabel{ord:oneplus} $1 \ordplus  \alpha = \alpha$
	%\alpha \approx \alpha \ordplus 1$
	\item\ollabel{ord:plusone} $\alpha \approx \alpha\ordplus 1$;
	એટલે $\alpha$ અને $\alpha\ordplus 1$માં
	સમાન સંખ્યાના ઘટકો છે
	\item\ollabel{ord:infinite} $\alpha$ ડેડેકિન્ડ-અનંત છે
	\end{enumerate}
\end{lem}
\noindent
આમ ક્રમવાચક સંખ્યા (અ)સીમિત હોવાનો અર્થ સમજવાની
પાંચ સાબિત કરી શકાય એવી સમકક્ષ રીતો મળી.

\begin{proof}
\emph{\olref{ord:notinomega} $\Rightarrow$ \olref{ord:omegaplus}.}
ત્રિવિકલ્પિતાથી.

\emph{\olref{ord:omegaplus} $\Rightarrow$ \olref{ord:oneplus}.}
$\alpha \geq \omega$ નક્કી કરો. અનંતાતીત અનુમાનથી
એવી લઘુતમ ક્રમવાચક સંખ્યા $\gamma$ મળે છે
(શક્ય છે $0$) કે કોઈ સીમા ક્રમવાચક સંખ્યા $\beta$
માટે $\alpha = \beta \ordplus \gamma$. હવે:
\[
	1 \ordplus \alpha =
	1 \ordplus (\beta \ordplus \gamma) =
	(1 \ordplus \beta) \ordplus \gamma =
	\supstrict_{\delta < \beta} (1 \ordplus  \delta) \ordplus  \gamma =
	\beta \ordplus  \gamma =
	\alpha.
\]
\emph{\olref{ord:oneplus} $\Rightarrow$ \olref{ord:plusone}.}
સ્પષ્ટપણે એવું !!a{bijection}
$f \colon (\alpha \disjointsum 1) \to (1
\disjointsum \alpha)$ છે. જો $1 \ordplus \alpha = \alpha$
હોય, તો એવી ક્રમ-એકરૂપતા
$g \colon (1 \disjointsum \alpha) \to \alpha$
છે. હવે $\comp{f}{g}$ વિચારો.

\emph{\olref{ord:plusone} $\Rightarrow$ \olref{ord:infinite}.}
જો $\alpha \approx \alpha \ordplus 1$ હોય, તો એવું
!!a{bijection} $f \colon
(\alpha \disjointsum 1) \to \alpha$ છે. દરેક
$\gamma < \alpha$ માટે $g(\gamma) = f(\gamma, 0)$
વ્યાખ્યાયિત કરો; આ !!{injection} બતાવે છે કે
$\alpha$ ડેડેકિન્ડ-અનંત છે, કારણ કે
$f(0,1) \in \alpha \setminus \ran{g}$.

\emph{\olref{ord:infinite} $\Rightarrow$ \olref{ord:notinomega}.}
આ \olref[z][infinity-again]{naturalnumbersarentinfinite} છે.
\end{proof}

\end{document}
